% ============================================================
% 50_typography_semantics.tex — Schriftmakros für Box-Typen & semantische Inhaltsmakros
% ============================================================

% Explizite bp-Groessen (wie Analysis) fuer praezise, font-unabhaengige Kontrolle
\newcommand{\ZSFfontChapter}{\sffamily\bfseries\fontsize{8.1bp}{9.1bp}\selectfont}
\newcommand{\ZSFfontSection}{\sffamily\bfseries\fontsize{8.1bp}{9.1bp}\selectfont}
\newcommand{\ZSFfontBoxTitle}{\sffamily\bfseries\fontsize{8.1bp}{9.1bp}\selectfont}
% chktex 6
\newcommand{\ZSFfontNote}{\sffamily\fontsize{6.7bp}{7.5bp}\itshape\selectfont}

% Tabellen-Hilfsmakros
\newcommand{\TableCellNote}[1]{{\ZSFfontNote #1}}
\newcommand{\TableTitleBreak}{\newline}
\newcommand{\TableTitle}[2][]{%
  \textbf{#2}%
  \if\relax\detokenize{#1}\relax\else\TableTitleBreak\TableCellNote{#1}\fi
}

% Semantische Tabellen-Fontgroessen
% (werden als optionales Argument an ZSFtable/ZSFtableFlat/ZSFtablePlain
%  uebergeben — siehe styles/20_tables.tex).
% Default = \ZSFfontTable. Fuer lange, dichte Tabellen (viele Zeilen,
% knapper Inhalt) \ZSFfontTableDense verwenden.
\newcommand{\ZSFfontTable}{\normalsize}
\newcommand{\ZSFfontTableDense}{\footnotesize}

% Semantische Inhaltsmakros (konsistente Skalierung und einfache Anpassung)
%
% \ZSFkeyword{...} — hebt zentrale Fachbegriffe im Fliesstext hervor.
% Verwendung:
%   * nur in \begin{runintext}, \begin{defbox}, \begin{warnbox} und
%     aehnlichen Text-Boxen, NIEMALS in Formeln, \formulanote oder
%     Ueberschriften (\SubsectionBar, \StartChapter).
%   * sparsam einsetzen (Richtwert: 1-3 pro runintext, 4-6 im
%     Kapitel-Intro, insgesamt ca. 10 pro Kapitel) -- es sollen echte
%     Fachbegriffe (Gesetze, Effekte, Bauteile, zentrale Groessen)
%     sein, keine generischen Substantive oder Verben.
% Rendering ist bewusst schlicht als \textbf{...}; Aenderungen zentral
% hier vornehmen (z.B. Farbe, Underline) -- nicht in Kapiteln.
\newcommand{\ZSFkeyword}[1]{\textbf{#1}}
\newcommand{\ZSFScriptRef}[1]{\hfill\textcolor{white!80}{\small (S.#1)}}
\newcommand{\ZSFconclusion}[1]{%
  \par\vspace{-\parskip}\vspace{\ZSFspaceS}%
  \noindent$\Rightarrow$ #1%
  \par\vspace{-\parskip}\vspace{\ZSFspaceS}%
}

% Tabellen-Styling (benötigt [table]{xcolor} und array)
\newcommand{\ZSFrowColor}{\rowcolor{\FormulaboxBackColor}}
\newcommand{\ZSFtableRuleColor}{black!70}
\newcommand{\SetZSFtableRuleColor}[1]{\renewcommand{\ZSFtableRuleColor}{#1}}
\newcommand{\ZSFtableRuleWidth}{0.5pt}
\newcommand{\SetZSFtableRuleWidth}[1]{\renewcommand{\ZSFtableRuleWidth}{#1}}
\newcommand{\ZSFbeginTableRuleStyle}{%
  \begingroup
  \arrayrulecolor{\ZSFtableRuleColor}%
  \setlength{\arrayrulewidth}{\ZSFtableRuleWidth}%
}
\newcommand{\ZSFendTableRuleStyle}{\endgroup}
